\documentclass{article}
\usepackage[margin=1in]{geometry}
\usepackage{amsmath, amsthm, booktabs, enumitem, listings}
\usepackage{forest}

\everymath{\displaystyle}
\newtheorem*{statement*}{Statement}

\title{Algorithms in C}
\author{}
\date{}

\begin{document}
\maketitle
\setcounter{section}{1}
\section{Principles of Algorithm Analysis}
\subsection{Implementation and Empirical Analysis}
\subsection{Analysis of Algorithms}
\subsubsection*{Exercises}
\begin{enumerate}[label=2.\arabic*]
    \setcounter{enumi}{2}
  \item We'll count the following operations: assignment, increment, and comparison. For each loop
    of $N$ iterations, there are $N$ loop body executions, $N$ increments, and $N+1$
    comparisons. For example, the innermost loop has $N$ increment operations
    (\texttt{count++}), $N$ increment operations (\texttt{k++}), and $N+1$ comparison operations
    (\texttt{k < N}). There are a total of 4 init assignment operations. Thus, the grand total
    number of operation is
    \[4 + \left(\left(N+N+(N+1)\right) \times N+N+(N+1)\right) \times N+N+(N+1)
    = 5 + 3N + 3N^2 + 3N^3.\]
\end{enumerate}

\subsection{Growth of Functions}
\begin{statement*}
  $\lceil\lg(N+1)\rceil$ is the number of bits in the binary representation of $N$.
\end{statement*}
\begin{proof}
  Suppose that $k$ is the number of bits requried to represent N. We have,
  \begin{align*}
    2^{k-1} \le & N & \le & 2^k-1 \\
    2^{k-1}+1 \le & N+1 & \le & 2^k \\
    2^{k-1} < 2^{k-1}+1 \le & N+1 & \le & 2^k \\
    2^{k-1} < & N+1 & \le & 2^k \\
    k-1 < & \lg(N+1) & \le & k. \\
  \end{align*}
  The last inequality is equivalent to $\lceil\lg(N+1)\rceil=k$.
\end{proof}
\end{document}
